\chapter{Experimental Design}
\label{sec:experimental_design}

All experiments were conducted using a strict patient-level 5-fold design (see Appendix \ref{app:cv_schema}), ensuring that no patient contributed tissue to both the training and the evaluation set of a fold, while keeping class balance into account. Since several patients supplied more than one tissue section, this is a stronger guarantee than separating individual slides. We evaluated the proposed framework under three complementary settings: Train-on-Real, Test-on-Real (TR/TR), Train-on-Synthetic, Test-on-Real (TS/TR), and Train-on-Synthetic+Real, Test-on-Real (TSR/TR) \cite{esteban2017real}. We also evaluated the proposed Classification by Reconstruction (CbR) strategy as a reconstruction-based generative inference approach.

\section{Classification Tasks}

To evaluate the impact of biological signal strength on classification performance, all models were evaluated on two distinct tasks:
\begin{enumerate}
    \item \emph{Binary classification (IDH-Status).} Differentiating the IDH-Mutant from the IDH-Wildtype entails a large metabolic shift that is theoretically easier to detect above spectrometer noise.
    \item \emph{6-class grading.} Differentiating all six distinct glioma sub-grades is a more challenging task, as the biochemical differences between sub-grades are more subtle and the signal may not capture them.
\end{enumerate}

\section{Evaluation Protocols}
\paragraph{TR/TR Evaluation.} 
We established the real-data baseline using the Train-on-Real, Test-on-Real (TR/TR) paradigm. The training splits consisted exclusively of real spectra from the training folds. We evaluated four downstream classifiers:
\begin{enumerate}
    \item \emph{RF + SVM:} A Random Forest feature selector (top 20 bands) followed by a Support Vector Machine with RBF kernel and balanced class weights \cite{apollonov2024raman}.
    
    \item \emph{Random Forest:} A standard ensemble of 100 trees with balanced class weights.
    
    \item \emph{XGBoost:} A gradient-boosted tree ensemble utilizing histogram-based tree methods optimized for high-dimensional data.
    
    \item \emph{1D-CNN:} A deep convolutional classifier trained with early stopping.
\end{enumerate}

\paragraph{TS/TR Evaluation.}
To evaluate whether the CVAE captured discriminative structure from the training data, we employed the Train-on-Synthetic, Test-on-Real (TS/TR) framework. For each fold, the CVAE was trained exclusively on the corresponding real training split. By sampling from the latent prior $z \sim \mathcal{N}(0, I)$, the decoder was then used to generate synthetic spectra conditioned by class (40,000 for binary tasks and 120,000 for 6-class tasks) in almost all experiments, except in the TSR/TR ratio ablation study (see Table \ref{tab:appendix_exp_401} in Appendix). The downstream classifiers were trained exclusively on synthetic spectra and evaluated on the corresponding held-out real validation fold.

\paragraph{TSR/TR Evaluation.}
We also evaluated a Train-on-Synthetic+Real, Test-on-Real (TSR/TR) framework in which synthetic spectra were combined with the real training data. This setting evaluates whether synthetic augmentation improves classification performance relative to the TR/TR baseline.

\paragraph{CbR Evaluation.}
Finally, the proposed Classification by Reconstruction (CbR) inference strategy was evaluated on the held-out real validation folds. Reconstruction-based predictions were compared directly against the TR/TR and TS/TR baselines.

\section{Training Details}
We present here the specific details of our training setup. We note that the ML models were not tuned, as the focus of this study is the effect of data augmentation, not tuning the best classifiers. Most of our experiments are ablation designs where we test the effect on the accuracy.
\begin{enumerate}
    \item \textit{Spectrum region crop.} In all our experiments, we cropped the spectra to the fingerprint region ($700 - 1800 \text{ cm}^{-1}$) \cite{Harris2023}, thus we selected 571 spectral bands out of the $N_\lambda = 1738$. 
    \item \textit{Maximum number of spectra cap per slide.} We randomly drew a sample of a maximum of 500 spectra from each slide, ensuring that the slides are equally represented. We performed this operation to mitigate the high variability in the number of spectra per slide due to varying slide sizes and varying retention rates due to quality (see Fig. \ref{fig:tumor_retention_rates}). We drew a fixed sample for the classifiers and opted for a dynamic resampling for the CVAE training.
    \item \textit{CVAE training data augmentation:} We performed dynamic data augmentation at train time using the following operations: 
    (a) \textit{within-class mixup:} we used linear combinations of spectra within the same class to expand the input space; $x = \lambda x_1 + (1-\lambda) x_2$, with $\lambda \sim U(0,\alpha)$, and we set $\alpha = 0.5$; 
    (b) \textit{spectral shift:} we shifted the spectra by a random integer in $[-\text{max\_shift},+\text{max\_shift}]$ to simulate Raman calibration drift; we set $\text{max\_shift}=3$ in our experiments.
\end{enumerate}


\section{Statistical Evaluation}
\label{sec:statistical_eval}

Two protocols in this work group spectra into clusters, and they deliberately use different
units. It is worth stating the distinction explicitly, since both are referred to as
``clustering'' of the data and they serve different purposes.

\emph{Data splitting is performed at the patient level.} No patient contributes tissue to both
the training and the evaluation side of any fold (Section~\ref{sec:cv_design}). The purpose here
is to guarantee that evaluation data is genuinely unseen, so the correct unit is the strictest
one available: a model must not be tested on a patient it has already been shown, regardless of
which tissue section that patient supplied. Every accuracy reported in this work is obtained
under this constraint.

\emph{Uncertainty estimation is performed at the slide level.} The bootstrap described below
resamples whole slides rather than whole patients. The purpose here is different: not to prevent
leakage, but to estimate how much a reported metric would vary if the study were repeated. Each
slide corresponds to one Raman acquisition on one tissue section, and the dominant sources of
variation between measurements, i.e., instrument state, focusing, calibration drift, and sample
preparation, act at the level of the acquisition rather than at the level of the patient. Since
the object under evaluation is a spectral generator, we regard the acquisition session as the
natural exchangeable cluster, and resampling slides simulates the variability that a fresh batch
of acquisitions would produce. This is a deliberate relaxation relative to the splitting
protocol: where a single patient contributed several sections, those sections are treated as
exchangeable during resampling, which we consider appropriate for characterising spectral
variability even though it would not be appropriate for constructing a test set.

Due to the hierarchical structure of the dataset, individual spectra originating from the same
biopsy sample are strongly correlated. Consequently, statistical evaluation at the
individual-spectrum level violates the independent and identically distributed (I.I.D.)
assumption and may artificially underestimate uncertainty. Conversely, estimating variance solely
across the five cross-validation folds yields limited statistical power.

To obtain more robust uncertainty estimates, we implemented a stratified cluster bootstrap
procedure at the slide level. Across 1,000 bootstrap iterations, slides were sampled with
replacement while preserving class proportions. This procedure preserves the correlation
structure within each slide while approximating variability across independent batches of biopsy
slides.

For pairwise comparisons between evaluation settings (TS/TR vs. TR/TR or TSR/TR vs. TR/TR), we computed bootstrap distributions of the corresponding accuracy differences:
\[
\Delta = \mathrm{Accuracy}_{A} - \mathrm{Accuracy}_{B}.
\]

Statistical significance was evaluated using bootstrap $p$-value estimates. Confidence Intervals (CI) were determined using the bootstrap distribution, and a result is considered significant when the 95\% CI around $\Delta$ excludes zero. The two sided $p$-value is defined as \[p = 2 \cdot \min\!\left(\frac{\# \{ i = \overline{1, N} : \Delta_i<0\}}{N},\frac{\#\{i = \overline{1, N} :\Delta_i>0\}}{N}\right),\] where $N = 1,000$. We note that no statistical correction is applied for multiple tests across the suite of experiments or classifiers; this is an important limitation and inflates the Type~I error rate.

